\answer{ex:09:1}{(a)~$P_\infty=0.2$, thời gian nhớ $20\,\text{s}$. (b)~$P_\infty=1$, thời gian nhớ $1\,\text{s}$. (c)~Thời gian nhớ tỷ lệ với biên độ nhiễu: sáu lần.}
\answer{ex:09:2}{$P(t)=P_\infty\coth(t\sqrt{q/R})$. (a)~$\frac12\sqrt{R/q}=5\,\text{s}$ --- ngắn bằng một nửa thời gian nhớ. (b)~$t=\frac12\ln21\,\sqrt{R/q}\approx15\,\text{s}$: \nt{ACET} phải giữ ở mức cao chừng ấy thời gian.}
\answer{ex:09:3}{Phương sai $qT/2+R/(2T)$; $T^*=\sqrt{R/q}$, phương sai $\sqrt{qR}=P_\infty$, như ở bộ lọc~\eqref{eq:kalman}; tăng $(\kappa+1/\kappa)/2$ lần: $1.25$ và $5.05$.}
\answer{ex:09:4}{$a>a_w=1-\theta/(mw)$: $0.15$ khi $w=0.041$, $0.22$ khi $w=0.045$.}
\answer{ex:09:5}{Lỗi xuất hiện khi $M\approx40$--$60$, khi $M=80$ --- $1.9$ nơ-ron thừa, khi $M=100$ --- tỷ lệ mẫu $0.86$; nhiễu từ các mẫu khác có phương sai $\approx(M-1)p(1-p)/N$, nên $M_{\max}\approx80$, $M_{\max}\propto N/p$.}
\answer{ex:09:6}{Ví dụ, $\eta(a)=\eta\min\bigl(1,\max(0,(a-0.3)/0.6)\bigr)$. Với $\eta a$, cả ba nơ-ron lạ đều lọt vào mẫu; với $\eta(a)$ --- không nơ-ron nào, việc ghi khi $a\ge0.9$ vẫn như cũ (tỷ lệ $1.00$).}
\answer{ex:09:7}{(a)~$40$ lần phát lại; với $\eta(a)$ ở bài tập~\ref{ex:09:6} --- không bao giờ. (b)~$200$ lần phát lại, $10$ đêm.}
